\section{Data and experimental setup}
\label{sec:experimental-setup}

\subsection{Datasets and preprocessing}
\label{sec:data}

We use the DBLP-Scholar and Abt-Buy datasets from the Magellan repository
\citep{konda2016magellan}. The two datasets differ considerably in both scale and
the form of their textual information. DBLP-Scholar contains bibliographic records,
whereas Abt-Buy contains product records collected from two online retailers.
\rev{Our evaluation reformulates these benchmark datasets as record-level candidate
retrieval and linkage problems, with FLR and MMR calculated over file-$A$ records and
their candidate sets. These error rates are therefore not directly comparable to the
pair-level precision, recall or F1 scores reported under the standard entity-matching
classification protocol in, for example, DeepMatcher and Ditto
\citep{mudgal2018deep,li2020deep}.}
Table \ref{tab:dataset-summary} summarises the data used in the experiments.

\begin{table}[ht]
\centering
\caption{Summary of the datasets used in the experiments.}
\begin{tabular}{lrrrr}
\toprule
Dataset & $n_A$ & $n_B$ & Gold pairs & Matchable records in $A$ \\
\midrule
DBLP-Scholar & 2,616 & 64,263 & 5,347 & 2,408 \\
Abt-Buy      & 1,081 & 1,092  & 1,097 & 1,081 \\
\bottomrule
\end{tabular}
\label{tab:dataset-summary}
\end{table}

For DBLP-Scholar, file $A$ is the DBLP file and file $B$ is the Scholar
file. The original files contain the fields title, authors, venue and year.
The gold-standard mapping contains 5,347 pairs involving 2,408 distinct
DBLP records and 5,218 distinct Scholar records. A DBLP record may therefore
have more than one valid Scholar partner. Of the 2,616 DBLP records, 2,408
have at least one gold-standard match.

For Abt-Buy, file $A$ is the Abt file and file $B$ is the Buy file. The
original Abt file contains the product name, description and price, while the
Buy file additionally contains the manufacturer. The gold-standard mapping
contains 1,097 pairs. All 1,081 Abt records have at least one known match in
the Buy file, and some records have more than one valid partner.

For semantic blocking, we use the title field for DBLP-Scholar and the
product-name field for Abt-Buy. These fields contain no missing values in
either dataset. The other common text fields remain available for the field-level comparisons
used by the traditional scoring models. The complete
gold-standard mappings are retained when determining whether a candidate
set contains a valid match and when evaluating a predicted link.

For constructing supervised match examples, one valid file-$B$ partner is
retained for each matchable file-$A$ record. Thus, there are 2,408 retained
positive pairs for DBLP-Scholar and 1,081 for Abt-Buy. Importantly, the
other known valid partners are not treated as non-matches. At evaluation,
a predicted link is considered correct if the selected file-$B$ record is
any known valid partner of the corresponding file-$A$ record.


\subsection{Training and candidate construction}
\label{sec:training-candidates}

The experimental unit is a file-$A$ record together with its candidate
file-$B$ records, rather than an individual candidate pair. This distinction
is important because all the pairs associated with a given file-$A$ record
must remain together when the data are divided into training and test sets.

For the main experiments, the file-$A$ records are randomly divided 50--50
into training and test sets. In DBLP-Scholar this gives 1,308 records in
each half. In Abt-Buy, 540 records are assigned to training and 541 to
testing. The primary results reported in Sections 3 and 4 use random seeds
42, 43 and 44 and are averaged over the three resulting splits. \rev{In Tables
\ref{tab:comparison}--\ref{tab:discrimination-gap}, numbers in parentheses are
Monte Carlo standard errors, calculated as the sample standard deviation of
the three seed-level error rates divided by $\sqrt{3}$. The FLR and MMR
are calculated separately for each split before averaging. With only three
replications, these standard errors describe observed split-to-split
variation and do not support precise interpretation of small method differences.}
Seeds 45--51 are used later where additional Monte Carlo replications are
required.

Candidate blocks are constructed according to Scheme-II in Section 1.3.
First, the pretrained bi-encoder \texttt{all-MiniLM-L6-v2} is used to
embed the blocking text of every record in files $A$ and $B$. For each
record $i\in A$, the records in $B$ are ranked by cosine similarity of
their embeddings, and the $k$ most similar records form $C_i$. We consider
primarily $k=5$ and $k=50$, while $k=1$ is also useful for examining the
blocking floor.

The same procedure is applied independently to every file-$A$ record.
Consequently, the same file-$B$ record may occur in the candidate blocks
of several file-$A$ records. This mirrors how semantic blocking would be
applied to records outside the labelled sample and avoids imposing an
artificial one-to-one restriction at the blocking stage.

The labelled candidate pairs are then constructed for supervised learning. Each retained
valid partner provides one positive example for its matchable file-$A$ record. This retained
positive is supplied for supervised fitting whether or not that particular partner occurs
among the retrieved top-$k$ candidates. The non-match examples are drawn from the
retrieved candidate block: candidate-derived pairs that are not known valid partners are
treated as negatives, while every additional gold-standard partner of the same file-$A$
record is excluded from the non-match class. Thus, one retained positive may be accompanied
by several retrieved hard non-matches rather than by only one negative example. At
evaluation, the candidate block is not augmented with the retained positive, so a record
cannot be linked correctly when none of its valid partners has been retrieved. The resulting
non-matches are therefore relatively difficult examples selected by semantic similarity
rather than arbitrary pairs drawn from $A\times B$.

For DBLP-Scholar, the training and test sets also contain file-$A$ records
that have no gold-standard partner. Their candidate pairs are all
non-matches. This is useful because such records occur in the application
of record linkage and a classification procedure should be able to leave
them unlinked. For Abt-Buy, every file-$A$ record has at least one known
valid partner.

For discriminative link classification, the candidate pairs belonging to
the training half are further divided 80--20 into fitting and validation
parts. The fitting part is used to estimate the model and the validation
part to choose the link-classification threshold by maximising the F1
score. The outer test records are not involved in this choice.


\subsection{Models and implementation}
\label{sec:models-implementation}

We consider three ways of progressing from the candidate pairs to link
classification, corresponding to the workflows introduced in Figure
\ref{fig:workflow}.

In the traditional workflow, field-level text comparisons are obtained
with the Jaro--Winkler metric. In the proposed workflow, these are replaced
by cosine similarities calculated from bi-encoder embeddings. Unless
otherwise stated, the bi-encoder is
\texttt{all-MiniLM-L6-v2} \citep{reimers2019sentence}. Section
\ref{sec:embedding-scoring} explains how these comparison outcomes enter
the subsequent scoring models.

Two traditional scoring models are used in the main comparison. The
generative model follows the MEC approach of \citet{lee2021maximum},
with the match and non-match densities of the continuous comparison
outcomes estimated by kernel density estimation. The discriminative model
is logistic regression. A radial-kernel SVM, XGBoost and a small
multilayer perceptron are considered later as a sensitivity analysis.

For direct classification from paired texts, we use the cross-encoder
\texttt{ms-marco-MiniLM-L-6-v2}. Both transformer models are first
considered in their pretrained form. Their fine-tuning is then investigated
in Section \ref{sec:fine-tuning-models}, where we give the corresponding
training settings and explain how the candidate blocks are treated before
and after fine-tuning.

For discriminative models, including the cross-encoder, a validation
subset of the training pairs is used to choose the threshold for declaring
a link by maximising the F1 score. The particular classification rules are
described together with the corresponding models below, since they form
part of the methodological comparison rather than merely an implementation
detail.

The objective of these experiments is not to identify the best possible
model or hyperparameter configuration for either benchmark dataset.
Rather, the common setup enables us to examine how embedding, scoring,
classification and fine-tuning can be combined within the proposed
workflow. \rev{Accordingly, the reported error rates are intended to illustrate the
behaviour of this workflow under a common experimental setup, rather than to
establish state-of-the-art entity-matching performance or a benchmark ranking
against published systems such as DeepMatcher \citep{mudgal2018deep}, Ditto
\citep{li2020deep}, or the method of \citet{peeters2021dual}.}


\subsection{Evaluation}
\label{sec:evaluation}

The unit of link classification is again a file-$A$ record. Given the
candidate block $C_i$, a scoring model ranks or scores its candidate
file-$B$ records. If a link is declared, let $j_i$ denote the selected
record. For models that use a classification threshold, no link is
declared when the best candidate does not satisfy that threshold. The
generative MEC calculations considered below instead select the
highest-ranked candidate according to the estimated density ratio.

A declared link is counted as correct whenever $(i,j_i)$ is one of the
known gold-standard pairs for record $i$. This is important for both
datasets because the gold-standard relation need not be one-to-one.
Selecting a valid partner other than the particular partner retained for
supervised fitting is therefore still a correct link.

The two error measures are the FLR and MMR defined in Equation (1).
It is useful to state their empirical calculation explicitly. Let $D$
be the number of links declared among the test file-$A$ records, let
$C$ be the number of those declarations that select a valid
gold-standard partner, and let $M$ be the number of matchable file-$A$
records in the test set. Then

\[
\widehat{\mathrm{FLR}}
=
1-\frac{C}{D},
\qquad
\widehat{\mathrm{MMR}}
=
1-\frac{C}{M}.
\]

Thus, an incorrect declared link contributes to both error measures:
it is a false link among the declarations and, for a matchable
file-$A$ record, it also means that the true match has not been
recovered. An undeclared matchable record contributes to the MMR but
not to the numerator or denominator of the FLR. For DBLP-Scholar,
a declared link for a non-matchable file-$A$ record is necessarily a
false link and therefore contributes to the FLR.

Blocking imposes an additional lower bound on the MMR. If none of the
valid partners of a matchable file-$A$ record occurs in $C_i$, no
scoring or classification model operating on that block can link the
record correctly. We refer to the proportion of such records as the
floor MMR and examine it explicitly in Section
\ref{sec:embedding-scoring}.

Unless otherwise stated, the tables in Sections 3 and 4 report the mean
FLR and MMR over the three primary random splits. In the later
investigation of training-set fractions and subsample error assessment,
seeds 42--51 provide ten Monte Carlo replications. There we additionally
report Monte Carlo standard errors, and, where relevant, relative Monte
Carlo standard errors, in order to distinguish variation across random
splits from the differences between modelling procedures.

\rev{The replication counts were not selected to attain a prespecified
hypothesis-testing power or precision target. They reflect the computational
cost of repeated fine-tuning and the descriptive purpose of the experiments.
Accordingly, the three-split comparisons in Tables
\ref{tab:comparison}--\ref{tab:discrimination-gap} are used to identify broad
patterns rather than to resolve small numerical differences; their Monte Carlo
standard errors are reported to show the scale of split-to-split variation. The
ten-replication studies in Section 4.3 provide a more precise assessment of
training-fraction sensitivity and subsample error estimation, but they are
likewise interpreted cautiously when the Monte Carlo standard error is of the
same order as the observed difference or bias.}
